% Section 4.3, from lectures/L04/README.md: the objectives, the evaluation
% questions, and the next lecture.

\section{Review}
\label{c4:sec:review}

\needspace{6\baselineskip}
\subsection*{What you should be able to do now}
After this chapter, you should:

\begin{itemize}
  \item Be able to choose the right kernel object for a synchronization or communication problem.
  \item Understand why a semaphore used as a lock allows unbounded priority inversion, and a mutex does not.
  \item Be able to pass data between threads with a message queue, and handle a full queue deliberately.
  \item Understand how Zephyr applications avoid dynamic memory, and what to use when they cannot.
  \item Be able to decouple an application’s parts with zbus channels, and choose between its kinds of
    observer.
  \item Be able to express a device’s behaviour as a hierarchical state machine with SMF.
\end{itemize}

\needspace{6\baselineskip}
\subsection*{Questions to test yourself}
You should be able to explain:

\begin{enumerate}
  \item What is the difference between a semaphore and a mutex, beyond the count?
  \item What is priority inversion, and why is the medium-priority thread the one that makes it unbounded?
  \item How does priority inheritance bound it, and what does it not solve?
  \item When should a producer use \code{K_NO_WAIT} on a full queue, and when a timeout?
  \item Why do Zephyr applications usually allocate kernel objects and buffers statically?
  \item What is the difference between a zbus listener, a subscriber and a message subscriber, and in which
    thread does each one run?
  \item Why can a zbus subscriber miss a message, and which observer cannot?
  \item What does a parent state in SMF save you from writing?
  \item On \code{native_sim}, why must a thread that simulates CPU work use \code{k_busy_wait()} rather than
    a loop?
\end{enumerate}

\needspace{6\baselineskip}
\subsection*{Looking ahead}
Chapter~\ref{c5:ch} is about writing a driver:

\begin{itemize}
  \item The device driver model: API structs, config and data, instances and init levels.
  \item A binding and a Kconfig option for a new part.
  \item Emulators: hardware that exists only in software, on \code{native_sim}’s I2C bus.
  \item Testing the driver with ztest, and running the tests with Twister.
\end{itemize}

You will write a driver for an I2C temperature sensor that does not exist.
